% GENERATED FILE. Edit canonical lessons, then run npm run generate:books.
% canonical-source-hash: fnv1a64:94ee59658d0a53db
% canonical-lessons: FR-C125-pompier, FR-C125-facteur, FR-C125-serveur, FR-C125-client, FR-C125-caisse

\chapter{Firefighter, Postman, Waiter, Customer, Till}
\label{ch:fr-firefighter-postman-waiter-customer-till}
\hlchaptermodality{125}

\begin{chapteropening}
\textbf{By the end of this chapter:} \emph{I can name a firefighter, a postman, a waiter, a customer, a till.}

Complete the last lesson of chapter 125: I can name a firefighter, a postman, a waiter, a customer, a till.
\end{chapteropening}

\section[le pompier]{le pompier — a firefighter}

\label{lesson:FR-C125-pompier}



\begin{quote}
Before the new one: say the French for a lawyer; an avocado, then the French for an engineer.
\end{quote}

\subsection*{You'll want to know: le pompier}
\textbf{le pompier} — \textquotedblleft{}a firefighter\textquotedblright{}.

\begin{cousinweb}
\textbf{pompier} — built from the word for \textquotedblleft{}a pump\textquotedblright{}, with \emph{-ier}: the one who works the pump.
\end{cousinweb}

\begin{grammarlens}[title={What you've built}]
One more word for this chapter.
\end{grammarlens}

\subsection*{Your turn}
\begin{itemize}
\raggedright
  \item \emph{Say it:} \emph{le pompier}
  \item \emph{Say it:} \emph{le pompier}, once more
  \item \emph{Say it:} say \emph{le pompier}
  \item \emph{Recall:} say the French for a lawyer; an avocado, then the French for an engineer, then say \emph{le pompier} again
\end{itemize}

\begin{tcolorbox}[breakable,colback=teal!4,colframe=teal!35!black,title={Before you move on}]
What does le pompier mean? (\textquotedblleft{}A firefighter\textquotedblright{}.) Say it once more.
\end{tcolorbox}
\section[le facteur]{le facteur — a postman, a postwoman}

\label{lesson:FR-C125-facteur}



\begin{quote}
Before the new one: say the French for an engineer, then the French for a firefighter.
\end{quote}

\subsection*{You'll want to know: le facteur}
\textbf{le facteur} — \textquotedblleft{}a postman, a postwoman\textquotedblright{}.

\begin{cousinweb}
\textbf{facteur} $\leftarrow$ Latin \emph{factor}, \textquotedblleft{}a maker, a doer\textquotedblright{}, from \emph{facere}, \textquotedblleft{}to make, to do\textquotedblright{} (EtymDB 2.1) — a learned borrowing; the agent who delivers became the postman.
\end{cousinweb}

\begin{grammarlens}[title={What you've built}]
One more word for this chapter.
\end{grammarlens}

\subsection*{Your turn}
\begin{itemize}
\raggedright
  \item \emph{Say it:} \emph{le facteur}
  \item \emph{Say it:} \emph{le facteur}, once more
  \item \emph{Say it:} say \emph{le facteur}
  \item \emph{Recall:} say the French for an engineer, then the French for a firefighter, then say \emph{le facteur} again
\end{itemize}

\begin{tcolorbox}[breakable,colback=teal!4,colframe=teal!35!black,title={Before you move on}]
What does le facteur mean? (\textquotedblleft{}A postman, a postwoman\textquotedblright{}.) Say it once more.
\end{tcolorbox}
\section[le serveur]{le serveur — a waiter}

\label{lesson:FR-C125-serveur}



\begin{quote}
Before the new one: say the French for a firefighter, then the French for a postman.
\end{quote}

\subsection*{You'll want to know: le serveur}
\textbf{le serveur} — \textquotedblleft{}a waiter\textquotedblright{}.

\begin{cousinweb}
\textbf{serveur} — built from the verb for \textquotedblleft{}to serve\textquotedblright{}, $\leftarrow$ Latin \emph{servire}, with \emph{-eur}, \textquotedblleft{}one who\textquotedblright{}. (EtymDB 2.1 records the verb.)
\end{cousinweb}

\begin{grammarlens}[title={What you've built}]
One more word for this chapter.
\end{grammarlens}

\subsection*{Your turn}
\begin{itemize}
\raggedright
  \item \emph{Say it:} \emph{le serveur}
  \item \emph{Say it:} \emph{le serveur}, once more
  \item \emph{Say it:} say \emph{le serveur}
  \item \emph{Recall:} say the French for a firefighter, then the French for a postman, then say \emph{le serveur} again
\end{itemize}

\begin{tcolorbox}[breakable,colback=teal!4,colframe=teal!35!black,title={Before you move on}]
What does le serveur mean? (\textquotedblleft{}A waiter\textquotedblright{}.) Say it once more.
\end{tcolorbox}
\section[le client]{le client — a customer}

\label{lesson:FR-C125-client}



\begin{quote}
Before the new one: say the French for a postman, then the French for a waiter.
\end{quote}

\subsection*{You'll want to know: le client}
\textbf{le client} — \textquotedblleft{}a customer\textquotedblright{}.

\begin{cousinweb}
\textbf{client} $\leftarrow$ Latin \emph{cliens}, \emph{clientis}, \textquotedblleft{}a dependant, a follower\textquotedblright{} — a learned borrowing. (EtymDB 2.1)
\end{cousinweb}

\begin{grammarlens}[title={What you've built}]
One more word for this chapter.
\end{grammarlens}

\subsection*{Your turn}
\begin{itemize}
\raggedright
  \item \emph{Say it:} \emph{le client}
  \item \emph{Say it:} \emph{le client}, once more
  \item \emph{Say it:} say \emph{le client}
  \item \emph{Recall:} say the French for a postman, then the French for a waiter, then say \emph{le client} again
\end{itemize}

\begin{tcolorbox}[breakable,colback=teal!4,colframe=teal!35!black,title={Before you move on}]
What does le client mean? (\textquotedblleft{}A customer\textquotedblright{}.) Say it once more.
\end{tcolorbox}
\section[la caisse]{la caisse — a till, a checkout}

\label{lesson:FR-C125-caisse}



\begin{quote}
Before the new one: say the French for a waiter, then the French for a customer.
\end{quote}

\subsection*{You'll want to know: la caisse}
\textbf{la caisse} — \textquotedblleft{}a till, a checkout\textquotedblright{}.

\begin{cousinweb}
\textbf{caisse} $\leftarrow$ Old Provençal \emph{caissa} $\leftarrow$ Latin \emph{capsa}, \textquotedblleft{}a box\textquotedblright{} — borrowed from the south: an inherited form would begin with \emph{ch-}. (EtymDB 2.1 records the Latin word.)
\end{cousinweb}

\begin{grammarlens}[title={What you've built}]
One more word for this chapter.
\end{grammarlens}

\subsection*{Your turn}
\begin{itemize}
\raggedright
  \item \emph{Say it:} \emph{la caisse}
  \item \emph{Say it:} \emph{la caisse}, once more
  \item \emph{Say it:} say \emph{la caisse}
  \item \emph{Recall:} say the French for a waiter, then the French for a customer, then say \emph{la caisse} again
\end{itemize}

\begin{tcolorbox}[breakable,colback=teal!4,colframe=teal!35!black,title={Before you move on}]
What does la caisse mean? (\textquotedblleft{}A till, a checkout\textquotedblright{}.) Say it once more.
\end{tcolorbox}
